% Bab 5: Low-Fidelity Prototype
\chapter{Low-Fidelity Prototype}
\label{bab:low-fidelity-prototype}
\section{Screen/State List}

Daftar layar sesuai implementasi di \texttt{frontend/app}.

\begin{longtable}{|L{0.092}|L{0.184}|L{0.507}|L{0.217}|}
\caption{Screen/State List}\\
\hline
\cellcolor[HTML]{D9EAF7}\textbf{ID} & \cellcolor[HTML]{D9EAF7}\textbf{Screen/State} & \cellcolor[HTML]{D9EAF7}\textbf{Tujuan} & \cellcolor[HTML]{D9EAF7}\textbf{Related Requirement} \\ \hline
\endfirsthead
\hline
\cellcolor[HTML]{D9EAF7}\textbf{ID} & \cellcolor[HTML]{D9EAF7}\textbf{Screen/State} & \cellcolor[HTML]{D9EAF7}\textbf{Tujuan} & \cellcolor[HTML]{D9EAF7}\textbf{Related Requirement} \\ \hline
\endhead
UI01 & Login \& Registrasi (\texttt{auth/login}, \texttt{auth/register}) & Masuk atau membuat akun lokal sebelum memakai aplikasi & F11 \\ \hline
UI02 & Onboarding: Welcome & Memperkenalkan aplikasi: jelajah kuliner, cerita budaya, culinary trail & F09 \\ \hline
UI03 & Onboarding: Izin Lokasi & Menjelaskan pemakaian lokasi lalu meminta izin sistem & F01, F09 \\ \hline
UI04 & Onboarding: Tanpa Lokasi & Menjelaskan fitur yang tidak tersedia jika izin lokasi ditolak & F01, F09 \\ \hline
UI05 & Onboarding: Izin Notifikasi & Meminta izin notifikasi untuk trigger dan pengingat & F02, F04 \\ \hline
UI06 & Onboarding: Preferensi & Memilih preferensi rasa, jenis kuliner, dan opsi hanya tenant halal & F07, F08, F12 \\ \hline
UI07 & Explore (Beranda/Peta) & Menampilkan posisi pengguna, batas kawasan, marker POI, status area, jarak ke kawasan, dan Mode Jelajah & F01, F02, F06, F09 \\ \hline
UI08 & POI Trigger Notification & Notifikasi entry/proximity dengan tombol Dengar / Lewati di lock screen & F02, F04 \\ \hline
UI09 & POI Detail (\texttt{poi/[id]}) & Menampilkan foto, menu dan harga, jam buka, status halal, audio narasi, transkrip, serta ringkasan dan daftar ulasan & F03, F04, F10 \\ \hline
UI10 & Search & Pencarian toleran typo dengan filter harga, rasa, jenis, dan status buka & F05 \\ \hline
UI11 & Navigasi (\texttt{navigate/[poiId]}) & Rute jalan kaki, jarak, panah kompas, panduan audio, dan deteksi kedatangan & F04, F06 \\ \hline
UI12 & Trail Planner & Input durasi, budget, dan preferensi rasa untuk membuat trail & F08 \\ \hline
UI13 & Trail Itinerary & Urutan tenant, estimasi biaya dan waktu, progres, serta Lewati Tenant & F07, F08 \\ \hline
UI14 & Trail Summary & Ringkasan trail: jumlah tenant, total belanja, durasi, dan jarak & F08 \\ \hline
UI15 & Paspor Kuliner (\texttt{passport}) & Trail tematik, stempel digital, progres, badge, riwayat ulasan, dan tombol Tambah Ulasan & F10 \\ \hline
UI16 & Profil \& Pengaturan & Data akun, ringkasan paspor, edit preferensi (termasuk halal), akses paspor, halaman pengaturan, dan keluar akun & F09, F11, F12 \\ \hline
UI17 & Form Ulasan (modal) & Memilih tenant, memberi rating 1--5, memilih menu yang dimakan, menulis ulasan, dan mengunggah foto & F10 \\ \hline
\end{longtable}

\section{Wireframe}

\textit{Masukkan low-fidelity wireframe untuk screen/state utama. Prototype perlu memperlihatkan perubahan kondisi ketika pengguna masih berada di luar area POI, mulai mendekati POI, menerima pemicu, dan kemudian mengakses informasi eksplorasi.}

\textit{[Tempel low-fidelity prototype di sini]}
